\section{Turing Machines}
\begin{definition}[Turing Machine]
  A TM $M=\left( Q,\Sigma,\Gamma,\delta,q_0,B,F \right) $, where:
  \begin{itemize}
    \item $Q$: A finite set of states
    \item $\Sigma$: A finite set of input symbols. $\Sigma\subseteq \Gamma$
    \item $\Gamma$: A finite set of tape symbols
    \item $\delta$: A transition function taking as argument a state and stack
      symbol and returns a state, stack symbol to write and direction to move.
      $\delta:Q\times \Gamma\to Q\times \Gamma\times \left\{ L,R \right\} $ 
    \item $q_0$: A start state, $q_0\in Q$ 
    \item $B$: A blank symbol. $B\in \Gamma-\Sigma$
    \item $F$: A set of final states. $F\in \mathcal{P}\left( Q \right) $
  \end{itemize}
  \begin{figure}[h!]
    \centering
    \includegraphics[width=0.5\textwidth]{tm}
  \end{figure}
  A TM consists of a finite control, which can be in any of a finite set of
  states, a infinite tape divided into cells, each holding one of a finite set
  of symbols.\\

  Initially, the input is written on the tape with blank symbols infinitely to
  the left and right of the input. The tape head is positioned on the leftmost
  cell holding an input symbol.\\

  The transition function $\delta$ encapsulates a move of the turing machine,
  consisting:
  \begin{enumerate}
    \item A change in state in the control.
    \item A replacement of the tape symbol in the read cell.
    \item A move to the left or right of the current cell.
  \end{enumerate}
\end{definition}

We have alternative notations for the descriptions of TMs, similarly to DFAs
and PDAs, the transition diagram and table.

\begin{figure}[h!]
  \centering
  \includegraphics[width=0.6\textwidth]{tm-table}
\end{figure}

We use the transition table to show the moves of the TM, depending on the
current state and tape symbols.

\begin{figure}[h!]
  \centering
  \includegraphics[width=0.5\textwidth]{tm-diag}
\end{figure}

The transition diagram of a TM consists of a set of nodes corresponding to the
states of the TM, within which each arc from state $q$ to $p$ is labelled by
some $X/YD$, where $X,Y\in \Gamma$ and $D\in \left\{ L,R \right\} $.
Additionally, we denote accepting states by double circles and the start state
with an arrow labelled `start' entering that state.

\begin{definition}[Instantaneous Descriptions]
  The \textbf{ID} of a TM, analogous to that of the PDA, is the string
  \[X_1X_2\cdots X_{i-1}q{X}_{i}X_{i+1}\cdots X_{n}.\] 
  where $q$ is the state of the TM, the tape head is scanning the symbol $X_i$
  and $X_1 X_2 \cdots X_{n}$ is the portion of the tape from the leftmost to
  the rightmost nonblank symbols, except if the head is past the ends, in which
  case we extend to show the blank symbols.
\end{definition}

\begin{notation}[$\vdash$]
  If $\delta\left( q,X_i \right) =\left( p,Y,L \right) $, we denote
  \[X_1X_2\cdots X_{i-1} q X_i X_{i+1}\cdots X_{n}\vdash X_1X_2\cdots X_{i-2} q X_{i-1} Y X_{i+1}\cdots X_{n}.\] 
  This represents one move of the TM.
\end{notation}

\begin{notation}[$\vdash^{*}$]
  As usual, we denote $\vdash^{*}$ to represent zero or more moves of the TM.
\end{notation}

\subsection{Turing Machines}
\subsubsection{Acceptance by a TM}
\begin{definition}[Language of a TM]
  The \textit{language} of a TM $M$ is denoted $L(M)$ and is defined by
  \[L(M)=\left\{ x|\exists q_f\in F, q_0x\vdash^{*}\alpha q_f \beta \right\}.\]
  That is, it is the set of strings that, given as input, would eventually
  transit $M$ into an accepting state.
\end{definition}

\begin{definition}[Halting]
  We say a TM halts if it enters a state $q$, scanning a tape symbol $X$, and there
  is no move defined. $\delta\left( q,X \right)$ is undefined.
\end{definition}

\begin{remark}
  By un-defining $\delta\left( q_f, X \right)$ for all final states $q_f$ and
  tape symbols $X$, without changing the language accepted, the TM halts if it
  accepts. We can hence assume that a TM always halts when it accepts a string.
\end{remark}

Since language acceptance is equivalent to function computation, we further
define the notion of function computation on a TM. Provided an input $x$, we
interpret the content of the tape of the TM if/after it halts as the output
$f(x)$.

\subsubsection{Models Equivalent to TMs}
We present certain modifications or improvements to TMs to aid in their
specifications. These modifications do not extend the class of languages
accepted, or extend the basic model of the TM, but only affect running times
and space complexities.

\paragraph{Storage in Finite Control} The finite control not only to represent
a position in the ``program'' of the Turing machine, but also to hold a finite
amount of data. The finite control becomes a tuple of the current state, and
the stored values and variables.

\paragraph{Subroutines} TMs can be thought of as a colection of
``subroutines'', sets of states that perform some useful process. Each
subroutine includes a start state, and a ``return'' state. As such, we may
define general subroutines that that allow us to ``jump'' to a specific return
state depending on the current start state.

\paragraph{Multiple Tapes} The tape of a TM composed of multiple ``tracks''.
Each symbol becomes a tuple of symbols and the tape alphabet consists of tuples
of the original tape alphabet. This TM is hence capable of simulating multiple
tapes.

\paragraph{Nondeterministic TMs} The transition function may instead return a
set of triples, from which the NTM can choose any of the the triples to be the
next move. The language accepted is hence defined as the strings such that,
given as input, there exist a sequence of moves that lead to an ID with an
accepting state.

\begin{theorem}
  If $M_N$ is a nondeterministic TM, there exists a deterministic TM  $M_D$
  such that $L(M_N)=L(M_D)$.
\end{theorem}

Intuitively, in $M_D$, we explore the IDs of $M_N$ ``breadth-first''. That is,
we find all IDs reachable in $0$ moves, then all IDs reachable in $1$ move,
then all IDs reachable in $2$ moves, as so on. This is done with a tape
operating as a queue of IDs to simulate a single move of.\\

We may further consider several restrictions of TMs, that do not affect the
capabilities of TMs.

\paragraph{Semi-Infinite Tapes}  A TM that is infinite only in one
direction (WLOG, to the right) can simulate a TM that is infinite in both
directions. Analogue to the simulation of multiple tapes, we allow each symbol
to be a duple of tape symbols, with each $i$-th cell from the origin be
$(x_{-i}, x_i)$, where $x_i$ is the $i$-th cell from the origin in the original
TM model.

\paragraph{No Blank Symbols Write} The TM is not allowed to write a blank
symbol. This, coupled with semi-infinite tapes, ensure that the tape is at all
times a prefix of nonblank symbols followed by an infinity of blanks. We may
achieve this by creating a new tape symbol $B'$ that effectively functions as
the blank symbol.

\paragraph{Two Stack PDAs} A PDA equipped with two stacks is as powerful as a
TM. We take one stack to simulate the tape symbols to the left of the head, and
the other to simulate the tape symbols to its right. Each move transits the
PDA's state, pops and reads from one stack, and pushes the corresponding output
to the other.


\subsection{Languages of a TM}
Following specifications of TMs and the languages that they accept, we have the
following for the language classes defined by TMs' capabilities. 
\subsubsection{Recursive Enumerable Language}
We first consider the languages that can be accepted by TMs.

\begin{definition}[Recursively Enumerable Languages]
  If $L$ is $L(M)$ for some TM $M$, then $L$ is a \textit{recursively enumerable
  (computably enumerable) language}.
\end{definition}

\paragraph{Codings}
We may define, for a string $x$ over $\left\{ 0,1 \right\} ^{*}$, its coding to be
\[i=\left( 1x \right) _2-1.\]
\begin{notation}
  $\omega_i$ denotes the string with code $i$.
\end{notation}

Further we may define a coding for a TM $M=\left( Q,\Sigma,\Gamma,\delta,q_1,B,\left\{ q_2 \right\}  \right)$ where
\begin{itemize}
  \item $Q=\left\{ q_1,q_2,\ldots,q_{n} \right\} $. If there are multiple
    accepting states, we may define new tape symbols and modify the transition function
    to compress into a single accepting state.
  \item $\Sigma=\left\{ 0,1 \right\} $, WLOG.
  \item $\Gamma=\left\{ X_1,X_2,\ldots,X_s \right\} $, where $X_1=0,X_2=1, X_3=B$.
  \item The directions $L=D_1$ and $R=D_2$.
\end{itemize}
We have, in this way, a coding for any transition
\[\delta\left( q_i,X_j \right) =\left( q_k,X_l,D_m \right)\qquad\to\qquad
0^{i}10^{j}10^{k}10^{l}10^{m}\quad i,j,k,l,m\in \mathbb{N}^*.\]
This allows for a coding for $M$ to be $C_1 11 C_2 11 C_3\ldots C_{n}$ where
$C_i$ is the coding of the $i$-th transition of the TM.
\begin{notation}
  $M_i$ denotes the TM with code $i$, and $W_i$ denotes the language accepted
  by the TM of code $i$.
\end{notation}

Hence, we have a function from the valid binary strings to the set of all TMs,
or the class of all RE languages.
\begin{remark}
  Not all binary strings is a valid code for a TM or a RE language. By
  convention, we take take $M_i$ to be the TM with one state and no
  transitions, and immediately halts on any input.
\end{remark}

\begin{example}[Diagonalisation Language]
  $L_d$ is not RE, where the diagonalisation language $L_d=\left\{ w_i |
  w_i\not\in L(M_i)  \right\} $ is the set of strings $\omega$ such that the TM
  whose code is $\omega$ does not accept $\omega$ as input.
\end{example}
\begin{proof}
  Suppose a TM $M_j$.
  \begin{enumerate}
    \item Case 1: $\omega_j\in L(M_j)$. $\omega_j\not\in L_d$ by definition.
    \item Case 2: $\omega_j\not\in L(M_j)$. $\omega_j\in L_d$ by definition.
  \end{enumerate}
  In all cases, $L(M_j)\neq L_d$. Since $M_j$ is an arbitrary TM, all TMs do
  not accept $L_d$, and it is not RE. \qed
\end{proof}

\begin{example}[Universal Language]
  The universal language $L_U=\left\{ \left( M,\omega \right) | M\text{ accepts
  }\omega \right\}$ is RE.
\end{example}
\begin{proof}
  Consider a TM which has four tapes: input tape (with codings of $M$ and
  $\omega$), tape to simulate $M$ (with $\omega$), tape for state of $M$,
  scratch tape. We may then simulate a step of $M$ by reading the state of $M$,
  the current cell of $M$, and finding the corresponding transition on the
  input tape. We hence transit the state of $M$, and write into the current
  cell of $M$. By construction, the universal TM accepts $L_U$ and $L_U$ is RE.
  \qed
\end{proof}

\subsubsection{Recursive Languages}
Within the RE languages, we may further define classes of language in which a
TM is able to tell us that a given string is not in the language. That is, it
halts for all possible input strings.

\begin{definition}[Recursive Languages]
  If $L$ is $L(M)$ for some TM $M$, and $M$ halts on all inputs $\omega\in \Sigma^{*}$,
  then $L$ is a \textit{recursive (decidable) language}.
\end{definition}

\begin{theorem}
  If $L$ is recursive, then $\bar{L}$ is recursive.
\end{theorem}
\begin{proof}
  Suppose $M=\left( Q,\Sigma,\Gamma,\delta,q_0,B,F \right) $ accepts $L$ and
  halts on all inputs. Consider a TM $\bar{M}$ such that
  \begin{itemize}
    \item The accepting states of $M$ are the nonaccepting states of $\bar{M}$
      with no transitions. $\bar{M}$ will halt without accepting.
    \item $\bar{M}$ has a new accepting state $r$ with no transitions.
      $\bar{M}$ will halt and accept.
    \item For each nonaccepting state $q\in Q-F$ and tape symbol $\alpha\in
      \Gamma$, if $\delta\left( q,\alpha \right) $ is not defined,
      $\bar{\delta}\left( q,\alpha \right) =r$.
    \item All other transitions from nonaccepting states are unchanged. $q\in
      Q-F$, $\alpha\in \Gamma$, $p\in Q$, if $\delta\left( q,\alpha \right)
      =p$, $\bar{\delta}\left( q,\alpha \right) =p$.
\end{itemize}
Since $M$ is guaranteed to halt, $\bar{M}$ is also guaranteed to halt, and
accepts exactly strings that $M$ does not accept.
\end{proof}

\begin{theorem}
  A language $L$ is recursive $\iff$ $L$ and $\bar{L}$ is RE.
\end{theorem}
\begin{proof}
  ($\implies$) If $L$ is recursive, $\bar{L}$ is recursive, and they both are RE.\\
  ($\impliedby$) Suppose TMs $M_1,M_2$ such that $L(M_1)=L$ and
  $L(M_2)=\bar{L}$. Consider a 2-tape TM $M$ that simulates $M_1$ and $M_2$ in
  parallel. If $\omega$ is given as input and is in $L$, $M_1$ halts and
  accepts, and $M$ accepts. If $\omega$ is not in $L$, $M_2$ halts and accepts,
  and $M$ rejects. In all cases, one of the TMs are guaranteed to halt and
  accept. Therefore, $M$ is guaranteed to halt and $L$ is recursive.
\end{proof}

With the two theorems, we conclude that given a language $L $, we have one of four cases:
\begin{itemize}
  \item Both $L$ and $\bar{L}$ are recursive.
  \item Neither $L$ and $\bar{L}$ is RE.
  \item $L$ is RE but not recursive, and $ \bar{L}$ is not RE.
  \item $\bar{L}$ is RE but not recursive, and ${L}$ is not RE.
\end{itemize}

\begin{example}[Complement of the Universal Language]
  The complement of the universal language $\bar{L_u}$ is not RE. Suppose that
  some TM $M$ accepts  $\bar{L_u}$. We can then construct a machine $M'$ for
  accepting $L_d$:\\
  $M'(\omega_i):$
  \begin{enumerate}
    \item Extract $i$ from $\omega_i$.
    \item Run $M$ on $\left( M_i,\omega_i \right) $.
    \item $M'$ accepts $\iff$ $M$ accepts.
  \end{enumerate}
  Then $L_d$ is RE. A contradiction. Hence $\bar{L_u}$ is not RE. \qed
\end{example}

\subsection{Undecidability}
The recursive languages corresponds to what is normally defined as an
algorithm, as a TM that is not guaranteed to halt may not give us enough
information to conclude that a string is not in the language. We are hence
unable to fully `solve' the problem.

\begin{theorem}[Church-Turing Thesis]
  Every effectively computable function can be computed by a Turing Machine.
\end{theorem}
The thesis is not a theorem as effectively calculable does not have a
mathematical definition. In any formalisation of computability, we have found
all models of computation to be equivalent to, or strictly weaker than, a
Turing Machine. 

\begin{definition}[Partial Recursive Functions]
  If $f$ is a function computed by some turing machine, then $f$ is
  \textit{partial recursive (partially computable)}.
\end{definition}

\begin{definition}[Recursive Functions]
  If $f$ is a function computed by some turing machine and $f$ is defined on
  all $\Sigma^{*}$, then $f$ is \textit{recursive (computable)}.
\end{definition}

Since function computation is equivalent to language acceptance, by the
Church-Turing Thesis, the class of recursive languages are intuitively the
`limits of computation' of traditional computers.

\subsubsection{Reductions}
\begin{definition}[Turing Reducibility]
  We say that a problem $P_1$ is reducible to $P_2$, $P_1\le _m P_2$, if there exists a recursive
  function $f$ such that
  \[x\in P_1\iff f(x)\in P_2.\] 
  That is, we can convert each instance of problem $P_1$ to some problem in
  $P_2$ that has the same answer. $P_2$ is as hard as or harder than $P_1$.
\end{definition}
\begin{figure}[h!]
  \centering
  \includegraphics[width=0.4\textwidth]{reduction}
\end{figure}
\begin{theorem}[Reductions]
  Suppose $P_1\le _m P_2$. Then
  \begin{enumerate}
    \item If $P_1$ is not recursive, $P_2$ is not recursive.
    \item If $P_1$ is not RE, $P_2$ is not RE.
    \item If $P_2$ is recursive, $P_1$ is recursive.
    \item If $P_2$ is RE, $P_1$ is RE.
  \end{enumerate}
\end{theorem}

To show that a problem $P$ is hard, we pick a problem $P'$ which is hard and
then show that $P'\le _m P$, by constructing a recursive function $f$ such that
$x\in P'\iff f(x)\in P$.

\begin{example}[TMs with an Nonempty Language]
  The set of encoded TMs with at least one string, $L_{ne}=\left\{ M|L\left( M
  \right) \neq \emptyset  \right\} $ is RE but not recursive.
\end{example}
\begin{proof}
  Consider a TM $M'$, such that:
  \begin{lstlisting}[mathescape]
  M'$(M)$:
    for $t=0$ to $\infty$:
      for $i=0$ to $t$:
        if M$(\omega_i)$ accepts within $t$ steps, M' accepts.
      endfor
    endfor
  endM'
  \end{lstlisting}
  The algorithm tries $\omega_1$ with $1$ step, then $\omega_1, \omega_2$ with
  $2$ steps, and so on. If $M$ accepts any string in the alphabet, it will halt
  and accept. Hence $L(M)=L_{ne}$ and $L_{ne}$ is RE.\\

  To show that $L_{ne}$ is not recursive, we consider the complement of
  $L_{ne}$, $L_e=\left\{ M|L(M)=\emptyset  \right\} $, and show that
  $\bar{L_U}\le _m L_e$, and hence $L_e$ is not RE. We want to find a recursive
  function that finds a TM $M'$ such that 
  \[L(M')=\emptyset \iff \omega\not\in L(M).\]
  Given input $\left( M,\omega \right) $, we define
  \begin{lstlisting}[mathescape]
  M'($x$):
    for $t=0$ to $\infty$:
      if M$(\omega)$ accepts within $t$ steps, M' accepts.
    endfor
  endM'
  \end{lstlisting}
  If $M$ accepts $\omega$, $M'$ accepts all strings and $L(M')\neq \emptyset $.
  If $M$ does not accept $\omega$, $M'$ does not accept any output.
  Since $\bar{L_U}$ is not RE, $L_{e}$ is not RE.\\

  $L_{ne}$ is RE but $L_e$ is not RE, hence they are not recursive.
\end{proof}


\subsubsection{Properties of RE Languages}
\begin{definition}[Properties]
  A \textit{property} of the RE languages is a subset of the $RE$ languages.
\end{definition}
\begin{definition}[Trivial Properties]
  A \textit{trivial property} is a property that is either empty or all RE
  languages. A property is nontrivial $\iff$ there exists a language that
  fulfils a property and a language that doesn't.
\end{definition}
\begin{definition}[Semantic Properties]
  A property $P$ is semantic if it pertains only the behaviour of the program,
  rather than its syntax. 
  \[L_1=L_2\implies L_1\in P\iff L_2\in P\]
\end{definition}

\begin{theorem}[Rice Theorem]
  All nontrivial semantic properties of the RE languages is undecideable.
\end{theorem}
\begin{proof}
  Suppose $\mathcal{P}$ a nontrivial property of RE languages. We show that
  $L_e\le L_\mathcal{P}$. That is, there exists a recursive function $f$ such
  that
  \[M\in L_e\iff f(M)\in L_\mathcal{P}.\]
  WLOG, assume $\emptyset $ satisfies $\mathcal{P}$ (otherwise, consider
  $\bar{\mathcal{P}}$). Let $L$ be a RE language that does not satisfy
  $\mathcal{P}$, and $M''$ the machine that accepts $L$. Given a TM $M$,
  we define $M'$:
  \begin{lstlisting}[mathescape]
M'$(x)$:
  for $t=0$ to $\infty$:
    for $i=0$ to $t$:
      if M$(\omega_i)$ accepts within $t$ steps $\wedge$ M''$(x)$ accepts within $t$ steps, M' accepts $x$. 
    endfor
  endfor
endM'
  \end{lstlisting}
  If $L(M)=\emptyset $, $L(M')=\emptyset $ that satisfies $\mathcal{P}$. If
  $L(M)\neq\emptyset $, $L(M')=L$ that does not satisfy $\mathcal{P}$. Since
  $L_e$ is not recursive, $L_\mathcal{P}$ is not recursive.
\end{proof}

\begin{remark}
  Rice's theorem does not imply that everything about a TM is undecidable. For
  instance, syntactic properties, such as the number of states of the TM, is
  decidable.
\end{remark}

\subsection{Undecidable Problems}
With the groundwork of undecidability and TMs laid out, we may move on to prove
the undecidability of problems involving beyond TMs and their languages. We
begin with the Post's Correspondence Problem, which involves strings rather
than TMs.

\subsubsection{Post's Correspondence Problem}
PCP consists two lists of strings, of equal length $k\in \mathbb{N}$, over an
alphabet $\Sigma$.
\[A=w_1,w_2,\ldots,w_k \qquad B=x_1,x_2,\ldots,x_k\]
Each $i\in \mathbb{N}_{\le k}$ specifies a \textit{corresponding pair} $\left(
w_i,x_i \right) $. A solution to PCP is a nontrivial sequence of integers
$i_1,i_2,\ldots,i_m$ taken as indices of corresponding pairs, that yield the
same string when concatenated.
\[w_{i_1}w_{i_2}\cdots w_{i_m}=x_{i_1}x_{i_2}\cdots x_{i_m}.\]
PCP is the decision problem for the existence of a solution.

\paragraph{``Modified'' PCP} To prove the undecidability of PCP, we first
consider a modified problem as an intermediate reduction. We define MPCP to be
an instance of PCP, with the additional constraint that the first corresponding
pair must also be the first pair in the solution. A solution similarly
comprises a sequence of integers $i_1,i_2,\ldots,i_m$ such that
\[w_1w_{i_1}w_{i_2}\cdots w_{i_m}=x_1x_{i_1}x_{i_2}\cdots x_{i_m}.\]

\begin{prop}
  MPCP is reducible to PCP.
\end{prop}
Suppose an instance of MPCP with lists $A=w_1,w_2,\ldots,w_k$ and
$B=x_1,x_2,\ldots,x_k$. WLOG, assume $*,\$\not\in\Sigma$. We construct a PCP
instance with $C=y_0,y_1,\ldots,y_{k+1}$ and $D=z_0,z_1,\ldots,z_{k+1}$ where:
\begin{itemize}
  \item For $i=1,2,\ldots,k$, $y_i:=w_i\cdot *$ and $z_i:=*\cdot x_i$
  \item $y_0:= *\cdot y_1$ and $z_0:=z_1$
  \item $y_{k+1}:=\$$ and $z_{k+1}:=*\$$
\end{itemize}
\begin{proof}
  ($\implies$) Suppose $i_1,i_2,\ldots,i_m$ is a solution to the MPCP instance.
  \[w_1w_{i_1}w_{i_2}\cdots w_{i_m}=x_1x_{i_1}x_{i_2}\cdots x_{i_m}.\]
  Applying the defined transformations, we have the equal strings
  \[*\cdot y_1y_{i_1}y_{i_2}\cdots y_{i_m}=z_1z_{i_1}z_{i_2}\cdots z_{i_m}\cdot *.\]
  Since $y_0=*\cdot y_1$, $z_0=z_1$, and $y_{k+1}=\$$, $z_{k+1}=*\$$, we have
  \[y_0y_{i_1}y_{i_2}\cdots y_{i_m}y_{k+1}=z_0z_{i_1}z_{i_2}\cdots z_{i_m} z_{k+1}\]
  and $0,i_1,i_2,\ldots i_m,k+1$ is a solution to the PCP instance.\\

  ($\impliedby$) Suppose a solution to the PCP instance. It must begin with
  index $0$ and end with index $k+1$ since they are the only corresponding
  pairs with the same start symbol and end symbol, respectively.
  \[y_0y_{i_1}y_{i_2}\cdots y_{i_m}y_{k+1}=z_0z_{i_1}z_{i_2}\cdots z_{i_m} z_{k+1}\]
  Removing the $*$s and the final $\$$, we have
  \[w_1w_{i_1}w_{i_2}\cdots w_{i_m}=x_1x_{i_1}x_{i_2}\cdots x_{i_m}\]
  and $1,i_1,i_2,\ldots,i_m$ is a solution to the MPCP instance.\qed
\end{proof}

\begin{prop}
  $L_U$ is reducible to MPCP.
\end{prop}
Suppose a pair $\left( M,\omega \right)$. WLOG, assume the TM never writes a
blank and the tape is semi-infinite (to the right). Each ID is hence of the
form $\alpha q\beta$ where $\alpha\in \Gamma^{+}, q\in Q, \beta\in \Gamma^{*}$.
For $M=\left( Q,\Sigma,\Gamma,\delta,q_0,B,F \right) $, we define the following
pairs in our MPCP instance
\begin{enumerate}
  \item $w_1=\#$ and $x_1=\#q_0w\#$ which must start the solution of MPCP to
    begin the simulation of $M$. 
  \item $w_2=x_2=\#$. For each $X\in \Gamma$, $w_i=x_i=X$.
  \item For all $q\in Q-F$, $p\in Q$, $X,Y,Z\in \Gamma$, 
    \begin{itemize}
      \item $w_i=qX$, $x_i=Yp$ if $\delta\left( q,X \right) =\left( p,Y,R \right) $
      \item $w_i=ZqX$, $x_i=pZY$ if $\delta\left( q,X \right) =\left( p,Y,L \right) $
      \item $w_i=q\#$, $x_i=Yp\#$ if $\delta\left( q,B \right) =\left( p,Y,R \right) $
      \item $w_i=Zq\#$, $x_i=pZY\#$ if $\delta\left( q,B \right) =\left( p,Y,L \right) $
    \end{itemize}
    Each of these pairs extend the $B$ string to add the next ID and extend the $A$ string
    to match the $B$ string.
  \item For all $q\in F$, $X,Y\in \Gamma$
    \begin{itemize}
      \item $w_i=XqY$, $x_i=q$
      \item $w_i=Xq$, $x_i=q$
      \item $w_i=qY$, $x_i=q$
    \end{itemize}
    These pairs allow us to consume all tape symbols after achieving an
    accepting state. This allows us to simplify to match the final step.
  \item $w_i=q\#\#$, $x_i=\#$. When the tape is emptied, we may use the final
    pair to complete the solution.
\end{enumerate}

Intuitively, the MPCP simulates the computation of $M$ on $\omega$, each
partial solution consisting strings that are prefixes of the sequence of IDs of
$M$, $\#\alpha_1\#\alpha_2\#\alpha_3\#\cdots$. The string from $B$ is one ID
ahead of that from $A$, unless $M$ enters an accepting state, where it would be
able to ``catch up'' to obtain a solution. Then $M$ accepts $\omega$ $\iff$
there is a solution to this instance of MPCP.\\

Since $L_U\le _m MPCP \le _m PCP$, by transitivity, since $L_U$ is undecidable,
PCP is undecidable.
\begin{theorem}
  PCP is undecidable.
\end{theorem}

\subsubsection{Ambiguous Grammars}
Using PCP, we may further prove the undecidability of other problems, particularly
that the ambiguity of a grammar is undecidable. Suppose an instance of PCP with lists
$A=w_1,w_2,\ldots,w_k$ and $B=x_1,x_2,\ldots,x_k$. WLOG, suppose
$a_1,a_2,\ldots,a_k\not\in \Sigma$. Consider the grammar, for $i\in \left[ 1,k
\right] $
\begin{gather*}
  S\to A|B\\
  A\to w_iAa_i| w_ia_i\\
  B\to x_iBa_i| x_ia_i
\end{gather*}
The grammar is ambiguous $\iff L(G_A)\cap L(G_B)\neq \emptyset \iff$ The
instance of PCP has a solution, where $G_A$ is the grammar of productions
involving $A$ and that of $B$ for $G_B$.

\begin{remark}
  $\bar{L\left( G_A \right) },\bar{L\left( G_B \right) }$ are CFLs.
\end{remark}

Using the undecidability of the intersection of $L(G_A)$ and $L(G_B)$, we may
prove the undecidability of various other decision problems.
\begin{example}[$L(G_1)=L(G_2)$]
  Take $L(G_1)=\bar{L(G_A)}\cup \bar{L(G_B)}=\bar{L(G_A)\cap L(G_B)}$ and
  $L(G_2)={\left( \Sigma\cup \left\{ a_1,a_2,\ldots \right\}  \right)}^{*}$.
\end{example}

\subsection{Unrestricted Grammars}
\begin{definition}[Unrestricted Grammar]
  A UG $G=\left( N,\Sigma,S,P \right) $, where:
  \begin{itemize}
    \item $N$: A finite set of non-terminals
    \item $\Sigma$: A finite set of terminals, $N \cap \Sigma=\emptyset $
    \item $S$: A start symbol, $S\in V$ 
    \item $P$: A finite set of productions
  \end{itemize}
  A production, contrary to CFGs, is of the form $\alpha\to\beta$, where
  $\alpha\in {\left( N \cup \Sigma \right)}^{*}N{\left( N\cup \Sigma
  \right)}^{*}$ and $\beta\in {\left( N \cup \Sigma \right)}^{*}$.
\end{definition}

\begin{theorem}
  The languages generated by unrestricted grammars are exactly the languages
  generated by TMs.
\end{theorem}

\begin{definition}[Context Sensitive Grammar]
  A CSG is a UG with the additional constraint that $\left| \alpha \right| \le
  \left| \beta \right| $.
\end{definition}
\begin{definition}[Context Sensitive Languages]
  If $L$ is $L(G)$ for some CSG $G$, then $L$ is a \textit{context-sensitive language}.
\end{definition}
\begin{example}[$\left\{ a^{n}b^{n}c^{n}|n\in \mathbb{N}^{*} \right\}$]
   \begin{gather*}
     S\to aSBC|aBC\\
     CB\to BC\\
     aB\to ab\\
     bB\to bb\\
     bC\to bc\\
     cC\to cc
   \end{gather*}
\end{example}

\begin{remark}
  The class of context sensitive languages are exactly $NSPACE\left( O\left( n
  \right)  \right) $. They are closed under intersection, union and concatenation.
\end{remark}
